% ============================================================================
% CAPACITY ANALYSIS OF THE BINARY REPRESENTATION
% ============================================================================
\subsection{Representational Capacity}
\label{sec:capacity_analysis}

This subsection analyzes the state space of the Value representation and
compares it with the capacity of conventional parameterized architectures.

\subsubsection{State Space of a Single Value}

A Value occupies $D = 8192$ bits of storage (128 words $\times$ 64
bits) and admits $2^{8192}$ possible binary configurations.  Of these,
the data region (words 0--59, yielding 3840 bits) contributes
$2^{3840}$ distinguishable data states, while the program region
(words 79--127, yielding 3136 bits encoding 98 instruction slots)
contributes $2^{3136}$ distinguishable programs.  Because data and
program regions are independently addressable, the joint state space of
a single frame is the Cartesian product of its regions---substantially
larger than either region alone.

\subsubsection{Combinatorial Interaction in Regions}

A Region holds up to 64 frames in its primary mixer and up to 256 in its
spill channel.  When two frames interact through a boolean operation, the
result depends on both frames' data and program contents.  For $N$ frames
in a Region, the number of ordered interaction pairs is $N(N-1)$, and
each pair can produce a distinct output depending on which boolean
operation is applied.  With 16 possible operations per instruction slot,
the reachable state space after $k$ interaction steps grows
combinatorially.

This combinatorial expansion is not merely theoretical: the
non-deterministic read order of the Region means that repeated
executions over the same initial set of frames may explore different
interaction sequences, providing a form of stochastic search over the
space of possible compositions.

\subsubsection{Comparison with Parameter-Count Scaling}
\label{sec:capacity_comparison}

Table~\ref{tab:capacity_comparison} contrasts the representational state
space of the Value architecture with that of conventional parameterized
models.

\begin{table}[h]
  \centering
  \caption{Memory footprint and representational state space.  The Value
           architecture achieves comparable information density per byte
           to parameterized models, but the structure of the state space
           differs: Values compose through boolean operations rather than
           matrix multiplications.}
  \label{tab:capacity_comparison}
  \begin{tabular}{lrrr}
    \toprule
    \textbf{Architecture}
      & \textbf{Memory}
      & \textbf{State Space}
      & \textbf{$\log_{10}$ States/KB} \\
    \midrule
    8192-dim float32 embedding
      & 32\,KB
      & $\sim 2^{262{,}144}$ (continuous)
      & $\sim 2{,}467$ \\
    GPT-2 (117M params, float16)
      & 234\,MB
      & $\sim 2^{1.87 \times 10^9}$
      & $\sim 2{,}406$ \\
    \addlinespace
    1K Values (mixed)
      & 1\,MB
      & $(2^{8192})^{1000}$
      & $\sim 2{,}467$ \\
    1M Values (mixed)
      & 1\,GB
      & $(2^{8192})^{10^6}$
      & $\sim 2{,}467$ \\
    \bottomrule
  \end{tabular}
\end{table}

The states-per-kilobyte ratio is comparable across architectures, as
information density is ultimately bounded by the entropy of the
underlying bits.  The difference lies in the structure of the state
space and the cost of navigating it:

\begin{itemize}[nosep]
  \item In a parameterized model, the state space is navigated through
        matrix multiplications and gradient updates over continuous
        parameters.  The cost per step scales with the number of
        parameters.
  \item In the Value architecture, the state space is navigated through
        boolean operations over binary frames.  Each operation executes
        in constant time per 64-bit word, independent of the number of
        stored Values.  Navigation is combinatorial rather than
        gradient-driven: the system explores states through composition
        of elementary boolean functions rather than through optimization
        of a differentiable objective.
\end{itemize}

\subsubsection{Implications}

The capacity analysis highlights a trade-off.  Parameterized models
achieve efficient navigation of their state space through learned
projections, at the cost of requiring gradient-based training and
floating-point arithmetic.  The Value architecture eliminates both
requirements but relies on combinatorial exploration through boolean
composition, which may require many interaction steps to reach a
desired state.  The experimental sections that follow provide empirical
evidence for what this exploration achieves in practice.
